\documentclass[10pt,a4paper]{amsart}
\usepackage{amsmath,amssymb,amsthm,graphicx}
% For commutative diagrams we use Xy-pic
\usepackage[all]{xy}

% We use 2cell for 2-commutative diagrams.
\xyoption{2cell}
% Theorem environments.
%
\theoremstyle{plain}
\newtheorem{theorem}[subsection]{Theorem}
\newtheorem{proposition}[subsection]{Proposition}
\newtheorem{lemma}[subsection]{Lemma}

\theoremstyle{definition}
\newtheorem{definition}[subsection]{Definition}
\newtheorem{example}[subsection]{Example}
\newtheorem{exercise}[subsection]{Exercise}
\newtheorem{situation}[subsection]{Situation}

\theoremstyle{remark}
\newtheorem{remark}[subsection]{Remark}
\newtheorem{remarks}[subsection]{Remarks}

\numberwithin{equation}{subsection}

% Macros
%
\def\lim{\mathop{\mathrm{lim}}\nolimits}
\def\colim{\mathop{\mathrm{colim}}\nolimits}
\def\Spec{\mathop{\mathrm{Spec}}}
\def\Hom{\mathop{\mathrm{Hom}}\nolimits}
\def\Ext{\mathop{\mathrm{Ext}}\nolimits}
\def\SheafHom{\mathop{\mathcal{H}\!\mathit{om}}\nolimits}
\def\SheafExt{\mathop{\mathcal{E}\!\mathit{xt}}\nolimits}
\def\Sch{\mathit{Sch}}
\def\Mor{\mathop{\mathrm{Mor}}\nolimits}
\def\Ob{\mathop{\mathrm{Ob}}\nolimits}
\def\Sh{\mathop{\mathit{Sh}}\nolimits}
\def\NL{\mathop{N\!L}\nolimits}
\def\CH{\mathop{\mathrm{CH}}\nolimits}
\def\proetale{{pro\text{-}\acute{e}tale}}
\def\etale{{\acute{e}tale}}
\def\QCoh{\mathit{QCoh}}
\def\Ker{\mathop{\mathrm{Ker}}}
\def\Im{\mathop{\mathrm{Im}}}
\def\Coker{\mathop{\mathrm{Coker}}}
\def\Coim{\mathop{\mathrm{Coim}}}

% Boxtimes
%
\DeclareMathSymbol{\boxtimes}{\mathbin}{AMSa}{"02}

%
% Macros for moduli stacks/spaces
%
\def\QCohstack{\mathcal{QC}\!\mathit{oh}}
\def\Cohstack{\mathcal{C}\!\mathit{oh}}
\def\Spacesstack{\mathcal{S}\!\mathit{paces}}
\def\Quotfunctor{\mathrm{Quot}}
\def\Hilbfunctor{\mathrm{Hilb}}
\def\Curvesstack{\mathcal{C}\!\mathit{urves}}
\def\Polarizedstack{\mathcal{P}\!\mathit{olarized}}
\def\Complexesstack{\mathcal{C}\!\mathit{omplexes}}
% \Pic is the operator that assigns to X its picard group, usage \Pic(X)
% \Picardstack_{X/B} denotes the Picard stack of X over B
% \Picardfunctor_{X/B} denotes the Picard functor of X over B
\def\Pic{\mathop{\mathrm{Pic}}\nolimits}
\def\Picardstack{\mathcal{P}\!\mathit{ic}}
\def\Picardfunctor{\mathrm{Pic}}
\def\Deformationcategory{\mathcal{D}\!\mathit{ef}}
\newenvironment{reference}{}{}
\newenvironment{slogan}{}{}
\newenvironment{history}{}{}
\title{Recurring matrices with nonnegative off-diagonal entries: the kernel is spanned by the vectors supported on the diagonal blocks}
\author{The Stacks Project Authors}
\date{stacks-project:models.tex:local-snapshot-20261012; GNU FDL 1.2 or later}
\begin{document}
\maketitle
\noindent\textit{Editorial note.} The chapter's candidate index records the proof body only, that is the lines strictly inside the proof environment; in this chapter file the line recorded as the first proof line is the first line after the opening marker. The retained proof span therefore starts one line earlier, at the opening marker itself, so that the assembled fragment is a complete proof environment. The recorded span 350-370 is the one that satisfies both begin and end markers, and it is the span used in proof\_spans and in the assembly body spans.\par\medskip
\noindent\textit{Editorial note.} The counted unit is the Stacks Project Lemma lemma-recurring-real at source lines 330-349 together with its proof at source lines 351-370, a straight adjacent pairing in the chapter file. The chapter is the upstream file models.tex, which is the file name recorded in the source record; the pool directory carries the longer name stacks-models and is not used as the chapter name anywhere in this card.\par\medskip
\noindent\textit{Editorial note.} The statement carries a labelled displayed equation, equation-ineq, at source lines 336-338, and the proof refers to it twice by number. That label belongs to the retained statement itself, so it resolves inside this card and no navigation pointer is needed.\par\medskip
\noindent\textit{Editorial note.} The statement and the proof contain no citation at all. The one citation mentioned in this card's editorial material is the reference environment of the neighbouring Lemma lemma-recurring of the same chapter, which this card does not count; it is described in the standard prerequisites so that the relation between the two statements is visible, and no bibliography entry has to be supplied for it.\par\medskip
\noindent\textit{Editorial note.} The retained preamble consists of the chapter preamble's native Xy-pic setup at lines 14-18, its theorem declarations and mathematical macros at lines 154-217, and the three no-op wrapper environments. The wrapper file defines reference, slogan and history as visible no-op environments; for this card that only matters for the definition of the reference environment, because the retained span does not use it.\par\medskip
\noindent\textit{Editorial note.} One asymmetry of the author's file is worth recording rather than repairing: in the neighbouring Lemma lemma-recurring the displayed inequality is not an equation with a label, whereas here the labelled equation equation-ineq is present and referred to by number, so nothing is missing from the retained span.\par\medskip
\noindent\textit{Editorial note.} The proof is the author's own complete argument: the rescaling step, the elementary inclusion, the maximal-coordinate equality chain and the induction on the number of nonzero coordinates. Nothing is omitted inside the proof environment, and the retained statement is a finished statement, so the card carries a complete unit.\par\medskip
\noindent\textit{Allowed standard prerequisites (editorial).}\begin{itemize}
\item Recurring matrices and the max-modulus argument. A real matrix A with nonnegative off-diagonal entries together with a positive vector m such that minus a of subscript i i times m of subscript i is at least the sum of a of subscript i j times m of subscript j over the other indices is what the chapter calls a recurring matrix for m; the classical diagonal-dominance argument of Taussky, quoted by the author in the neighbouring lemma lemma-recurring of this same chapter, gives the nonsingularity statement that the present lemma refines. The author's own definitions are retained in the span of the counted lemma.. The counted lemma is the structural refinement of the diagonal-dominance criterion for recurring matrices, and its proof is the same maximal-coordinate argument carried through with equality analysis.
\end{itemize}\medskip
\begin{lemma}
\label{lemma-recurring-real}
Let $A = (a_{ij})$ be a real $n \times n$ matrix with
$a_{ij} \geq 0$ for $i \not = j$. Let $m = (m_1, \ldots, m_n)$ be a real
vector with $m_i > 0$. For $I \subset \{1, \ldots, n\}$ let
$x_I \in \mathbf{R}^n$
be the vector whose $i$th coordinate is $m_i$ if $i \in I$
and $0$ otherwise. If
\begin{equation}
\label{equation-ineq}
-a_{ii}m_i \geq \sum\nolimits_{j \not = i} a_{ij}m_j
\end{equation}
for each $i$, then $\Ker(A)$ is the vector space
spanned by the vectors $x_I$ such that
\begin{enumerate}
\item $a_{ij} = 0$ for $i \in I$, $j \not \in I$, and
\item equality holds in (\ref{equation-ineq}) for $i \in I$.
\end{enumerate}
\end{lemma}


\begin{proof}
After replacing $a_{ij}$ by $a_{ij}m_j$ we may assume $m_i = 1$ for all $i$.
If $I \subset \{1, \ldots, n\}$ such that (1) and (2) are true,
then a simple computation shows that $x_I$ is in the kernel of $A$.
Conversely, let $x = (x_1, \ldots, x_n) \in \mathbf{R}^n$ be a
nonzero vector in the kernel of $A$. We will show by induction
on the number of nonzero coordinates of $x$ that $x$ is in the
span of the vectors $x_I$ satisfying (1) and (2). Let
$I \subset \{1, \ldots, n\}$ be the set of indices $r$ with $|x_r|$ maximal.
For $r \in I$ we have
$$
|a_{rr} x_r| = |\sum\nolimits_{k \not = r} a_{rk}x_k| \leq
\sum\nolimits_{k \not = r} a_{rk}|x_k| \leq
|x_r| \sum\nolimits_{k \not = r} a_{rk} \leq |a_{rr}||x_r|
$$
Thus equality holds everywhere. In particular, we see that
$a_{rk} = 0$ if $r \in I$, $k \not \in I$ and equality holds
in (\ref{equation-ineq}) for $r \in I$. Then we see that we
can subtract a suitable multiple of $x_I$ from $x$ to decrease
the number of nonzero coordinates.
\end{proof}

\end{document}
